% ID: SEL-D0038-OB-E04
% ESTADO: cotejada. Segunda pasada visual y borrador compilado revisados.
% Fuente: Selectividad/matematicas_aplicadas/Ua3MATEMÁTICAS APLICADAS examen.pdf
% SHA-256: 4440c15a31a77ee9e052a0668923508e117795b369d569bc02d24b7937366f7a
% Página: 2; opción B; ejercicio original 4.

La concentración de arsénico en los moluscos de una zona costera sigue una ley Normal con desviación típica 6 mg/kg. Para verificar la calidad de estos moluscos se toma una muestra aleatoria de tamaño 36 para contrastar si la media poblacional no supera el límite máximo de 80 mg/kg permitido por la normativa sanitaria $(H_0:\mu\leq80)$.
\begin{enumerate}[label=\alph*)]
\item \textbf{(1.5 puntos)} Determine la región crítica de este contraste a un nivel de significación del 5\%.
\item \textbf{(1 punto)} \textquestiondown{}Debe rechazarse esta hipótesis nula, al nivel del 5\%, si en esa muestra de 36 moluscos se encuentra una concentración media de arsénico de 82 mg/kg?
\end{enumerate}
